\section{What federation can and cannot bias}\label{ch:sums}

Before locating the loss, one fact narrows the search. It is also the idea on which the whole method rests.

\paragraph{A client split is a pure label shift} The Dirichlet split of \cref{ch:problem} decides \emph{how many}
images of each class a client gets, but within a class it picks the images uniformly at random. So the class-9 images
of client 7 are a random subset of all class-9 images, and the same is true of every client and class. In symbols, the
distribution of an image given its class is the same on every client as on the pooled data,
\[
P_k(x\mid y)=P(x\mid y)\quad\text{for every client $k$ that holds class $y$.}
\]
Only the class proportions $P_k(y)$ differ. The extractor is shared, so the same holds for the features $h=\phi(x)$.

The consequence: any quantity that describes one class, such as the mean feature of class 9, its spread, or the best
weight vector for class 9 in a last layer, has the \emph{same target} on every client that holds the class and on the
pooled data. Federation cannot move that target. It can only make the estimate noisier, because each client sees fewer
samples, or slower to learn. So when a class-level quantity ends up \emph{biased} in a federated method, something in
the method's estimation procedure must have caused it.

\paragraph{Sums are exact} Many class-level quantities can be computed from sums. The mean feature of class $c$ is
\[
\mu_c=\frac{\sum_k S^{(1)}_{kc}}{\sum_k n_{kc}},\qquad S^{(1)}_{kc}=\sum_{i\in D_{kc}}h_i,
\]
where $D_{kc}$ is client $k$'s set of class-$c$ images and $n_{kc}$ its size. If every client sends its count $n_{kc}$
and its sum $S^{(1)}_{kc}$, the server obtains $\mu_c$ exactly, whatever the split. For class 9 of the running
example, client 7 sends 63 and the sum of its 63 features, client 9 sends 14 and the sum of its 14, and so on; the
server divides the total sum by the total count, 100, and gets the same mean it would get from the 100 images in one
place. The paper states this as a lemma.

\begin{lemma*}[partition invariance]
An estimator that is a function of a sum $S=\sum_k\sum_{i\in D_k}T(x_i,y_i)$, with $T$ fixed in advance, takes the same
value for every partition of the pooled data, and equals the centralised estimator.
\end{lemma*}

The proof is that addition does not depend on the order or grouping of its terms. The lemma is simple, and the paper
says that it is known for closed-form federated classifiers (Fed3R, AFL, FedCOF). Its use here is that it tells which
quantities a federated method can compute \emph{without} the bias of \cref{sec:s-routes}: those that are functions of
per-client sums. The class spread is one. With $S^{(2)}_{kc}=\sum_{i\in D_{kc}}\|h_i\|^2$,
\[
V_c=\frac{\sum_k S^{(2)}_{kc}}{n_c}-\|\mu_c\|^2=\E\big[\|h-\mu_c\|^2\,\big|\,y=c\big],
\]
the mean squared distance of a class-$c$ feature from the class mean, is also exact.

\paragraph{Secure aggregation} Sending per-client sums to the server would reveal each client's class counts and
feature sums. \emph{Secure aggregation} is a cryptographic protocol that lets the server learn the sum over clients of
vectors and nothing else about any single vector. The basic idea: every pair of clients agrees on a random mask; one
adds it to its vector, the other subtracts it; each masked vector looks random to the server, but in the sum the masks
cancel. The paper uses Flower's SecAgg+, a variant that tolerates clients dropping out.

The protocol works on integers, so each vector is divided by a public bound and quantised to 30 bits with unbiased
stochastic rounding. The error is at most $1.5\times10^{-8}$ of the bound per client and coordinate. Values beyond
eight times the bound would be clipped, and with fewer than three participants the protocol cannot run and the sum is
computed in clear. The paper calls the statistics \emph{exact} in this sense: exact up to this quantisation. The
bounds come from a small validation set that the server holds (200 images per class on CIFAR-10, held out from the
training data), never from the clients' data.

What secure aggregation does not hide is the sum itself. If a class is held by a single client, the federation-level
sum of that class is that client's sum. \Cref{ch:method} returns to what the released sums reveal.

\paragraph{How bias still enters}\label{sec:s-routes} If the targets cannot be biased by the split, a federated method can
still bias its estimate of them in four ways, which the paper examines.
\begin{enumerate}
\item \emph{The optimiser.} An optimiser can treat clients that lack a class differently from clients that hold it,
and averaging can then pull a class parameter away from its target.
\item \emph{Parameters tied to one client's prior.} A last-layer bias trained on one client's data learns that
client's class proportions, which are not the federation's.
\item \emph{What the synthetic data teach.} Training on generated features teaches whatever those features carry. If
they do not carry their label, the training teaches something else.
\item \emph{Sample size.} Nonlinear estimates, such as a spread or an inverse covariance matrix, are biased in small
samples. Class 9 has 100 images and features of dimension 2560: its covariance matrix cannot even be inverted.
\end{enumerate}
The paper does not claim that the list is complete. For example, drift of the shared extractor in Stage I, where the
extractor is trained by FedAvg on skewed data, is not covered.

Which of these routes actually hurt GeFL-F?
